\documentclass[11pt]{amsart}
\usepackage[margin=1.15in]{geometry}
\usepackage{amsmath,amssymb,amsthm,mathtools}
\usepackage{microtype}
\usepackage[hidelinks]{hyperref}
\usepackage{booktabs}
\usepackage{alltt}
\emergencystretch=3em

\newtheorem{theorem}{Theorem}
\newtheorem{proposition}[theorem]{Proposition}
\newtheorem{lemma}[theorem]{Lemma}
\theoremstyle{remark}
\newtheorem{remark}[theorem]{Remark}

\newcommand{\Lam}{\Lambda}
\newcommand{\ord}{\operatorname{ord}}
\newcommand{\li}{\operatorname{li}}
\newcommand{\eps}{\varepsilon}
\newcommand{\Q}{\mathbb{Q}}
\newcommand{\Z}{\mathbb{Z}}
\newcommand{\C}{\mathbb{C}}
\renewcommand{\Re}{\operatorname{Re}}
\renewcommand{\Im}{\operatorname{Im}}
\newcommand{\lean}[1]{\texttt{#1}}
\newcommand{\repo}[1]{\href{https://github.com/senamakel/#1}{\texttt{#1}}}

\title[Seventeen consequences of RH in Lean]{Seventeen consequences of the Riemann hypothesis,\\ formalized in Lean by agents with persistent memory}
\author[S. Enamakel]{Steven Enamakel}
\email{enamakel@tinyhumans.ai}
\author[C. Buchanan]{Christel Buchanan}
\date{November 2026}

\begin{document}
\begin{abstract}
On 20 November 2026 Levent Alp\"oge released a Lean proof of the Riemann hypothesis and of its generalization to Dirichlet $L$-functions. We describe seventeen Lean~4 developments that combine that theorem with classical conditional results. Six prove named conjectures: Erd\H{o}s Problem 985, the Lindel\"of hypothesis, the density hypothesis, Grosswald's conjecture, the completeness of Euler's list of 65 idoneal numbers, and the extended Riemann hypothesis for abelian number fields. Five prove statements known to be equivalent to RH: the inequalities of Robin and Nicolas, Li's positivity, Speiser's theorem on $\zeta'$, and the Mertens bound $M(x)=O(x^{1/2+\eps})$. Six prove further bounds, most with the constants of the original papers, including part (iii) of Erd\H{o}s Problem 770 and Selberg's bounds toward Problem 233. Each development proves the implication from RH or GRH separately from RH itself, and checks the remaining finite ranges with programs whose soundness is proved. We give the shared analytic formulas, the main step of each proof, and the method: an atlas of 2{,}196 problems, a self-improving loop of agents whose only judge is Lean, and a source-linked memory used to find the targets.
\end{abstract}
\maketitle

\section{Introduction}

Number theory has many theorems of the form ``if the Riemann hypothesis holds, then $X$''. They were proved over the last century by Littlewood, Cram\'er, Selberg, Schoenfeld, Weinberger and others, and each was waiting for its hypothesis. Alp\"oge's proof \cite{alpoge} supplied it. The Lean FRO and the Mathlib Initiative certified the formal proof the next day \cite{leanfro}, and independent kernels accept its export \cite{sutherland}. The statements we use from it are
\begin{align}
\text{(RH)}\quad &\zeta(s)=0,\ s\notin\{-2,-4,\dots\}\ \Longrightarrow\ \Re s=\tfrac12, \label{eq:rh}\\
\text{(GRH)}\quad &\Re s>\tfrac12,\ (\chi\ne\chi_0\text{ or }s\ne1)\ \Longrightarrow\ L(s,\chi)\ne0, \label{eq:grh}
\end{align}
the second for every modulus $N\ge1$ and every Dirichlet character $\chi$ modulo $N$.

A proof of RH does not by itself produce a formal proof of its consequences. Someone has to find which conjectures depend on exactly \eqref{eq:rh} or \eqref{eq:grh}, formalize the conditional argument with explicit constants, and check the finite range that the analysis leaves over. Between 22 and 25 November 2026 we did this for seventeen targets, with agents run by OpenHuman \cite{openhuman} inside a loop called RSI Ralph \cite{rsiralph}, and with Living Brain \cite{livingbrain} as memory. The preparation began in August.

\subsection*{What is proved} Table~\ref{tab:main} lists the targets. All of them live in \url{https://github.com/ChatAndBuild/conjectures}, one repository per target. Rows 1 to 6 are named conjectures that were open before the proof of RH. Rows 7 to 11 are equivalent to RH, so proving them restates RH in another form; we formalize the direction from RH. Rows 12 to 17 are theorems that the literature proves under RH or GRH.

\subsection*{Credit} The proof of RH and GRH is Alp\"oge's. The conditional theorems belong to the authors cited in each section. Our contribution is the formal proof of each implication, the explicit constants and lemmas that the papers leave to the reader, the verified computations, and the resulting theorems without hypotheses. One statement, the bound in Theorem~\ref{thm:770} valid for every $n\ge2$, we did not find in the literature in that explicit form. Problems 770 (parts (i) and (ii)) and 233 remain open.

\begin{table}[ht]
\small
\setlength{\tabcolsep}{4pt}
\begin{tabular}{@{}rllll@{}}
\toprule
 & Repository & Result & Input & Conditional theorem\\
\midrule
1 & \repo{erdos-985} & prime primitive root $q<p$ & GRH & McGown et al.\ \cite{mtt}\\
2 & \repo{lindelof-rh} & Lindel\"of hypothesis & RH, GRH & Littlewood \cite{littlewood}\\
3 & \repo{density-rh} & density hypothesis & RH & classical\\
4 & \repo{grosswald-grh} & Grosswald's conjecture & GRH & McGown et al.\ \cite{mtt}\\
5 & \repo{weinberger-grh} & 65 idoneal numbers & GRH & Weinberger \cite{weinberger}\\
6 & \repo{abelian-erh} & ERH, abelian fields & GRH & see \cite{washington}\\
\midrule
7 & \repo{robin-rh} & Robin's inequality & RH & Robin \cite{robin}\\
8 & \repo{nicolas-rh} & Nicolas' inequality & RH & Nicolas \cite{nicolas}\\
9 & \repo{li-criterion-rh} & Li's positivity & RH & Li \cite{li}\\
10 & \repo{speiser-rh} & Speiser's theorem & RH & Speiser \cite{speiser}\\
11 & \repo{mertens-rh} & Mertens bound & RH & Littlewood \cite{littlewood}\\
\midrule
12 & \repo{lls-grh} & least non-residues, $L(1,\chi)$ & GRH & Lamzouri et al.\ \cite{lls}\\
13 & \repo{bach-sorenson-grh} & least prime $\equiv a\ (q)$ & GRH & Bach, Sorenson \cite{bachsorenson}\\
14 & \repo{schoenfeld-rh} & $\psi$, $\theta$, $\pi$ with $1/8\pi$ & RH & Schoenfeld \cite{schoenfeld}\\
15 & \repo{cramer-rh} & prime gaps, constant $\frac{22}{25}$ & RH & Carneiro et al.\ \cite{cms}\\
16 & \repo{erdos-770} & Erd\H{o}s 770(iii) & GRH & see \S\ref{sec:further}\\
17 & \repo{erdos-233} & bounds toward Erd\H{o}s 233 & RH & Selberg \cite{selberg}\\
\bottomrule
\end{tabular}
\medskip
\caption{The seventeen developments.}\label{tab:main}
\end{table}

\section{The formal setup}\label{sec:setup}

Every development has two layers. The main library depends on Mathlib and on the other repositories here, never on the proof of RH. Its headline theorem takes RH or GRH as an explicit hypothesis. A short file in a directory \lean{Unconditional/} then imports \lean{riemann-hypothesis} \cite{rhrepo} and supplies that hypothesis. For Erd\H{o}s 985 the whole file is
\begin{small}\begin{alltt}
theorem erdos_985 : Erdos985.Erdos985Statement :=
  Erdos985.erdos_985_of_GRH LFunction_ne_zero_of_half_lt_re
\end{alltt}\end{small}
The hypotheses are copied word for word from that repository, and both sides use Lean v4.36.0 and the same Mathlib revision, so the two definitions are the same term. If the main library says $\mathrm{RH}\Rightarrow X$, that implication stands whatever happens to the second layer.

Target statements sit in a file that is fixed before the proof starts. Where Google DeepMind's \emph{formal-conjectures} \cite{formalconj} has the statement, we copy it; otherwise we transcribe it from the paper. Fixing the statement matters because the easiest way to finish a hard proof is to weaken what it claims.

\subsection*{Finite ranges} An explicit proof usually works only above a threshold. Below it we run a checker, written in Lean, with a Lean proof that acceptance implies the theorem on that range. Large runs are evaluated by \lean{native\_decide}, which trusts the Lean compiler for the evaluation and records one axiom per call. Table~\ref{tab:compute} gives the counts; everything else depends only on Lean's axioms \lean{propext}, \lean{Classical.choice} and \lean{Quot.sound}. No development contains a \lean{sorry}.

\begin{table}[ht]
\small
\begin{tabular}{@{}llr@{}}
\toprule
Repository & Range checked by computation & \lean{native\_decide}\\
\midrule
\repo{erdos-985} & odd primes $p<10^8$ & 20\\
\repo{grosswald-grh} & all $409<p<10^{13}$; sieve tree above & 2950\\
\repo{weinberger-grh} & $n<10^8$; the 65 numbers & 11\\
\repo{robin-rh} & $\log n\le e^{30}$, by sieving to about $1.07\cdot10^{13}$ & 624\\
\repo{nicolas-rh} & primes $p_k\le10^5$ & 1\\
\repo{speiser-rh} & $\frac52\le|\Im s|<7$ & 2\\
\repo{lls-grh}, \repo{erdos-770} & moduli for Bach's constant $2$ & 15\\
\repo{bach-sorenson-grh} & $q<2^{15}$ & 16\\
\repo{schoenfeld-rh} & $\psi$ on $[1500,10^{10}]$; three small ranges & 24\\
\bottomrule
\end{tabular}
\medskip
\caption{Verified computations behind the headline theorems. The other seven use no \lean{native\_decide}; some of them run smaller checks in the kernel. In \lean{lls-grh} only the theorems with Bach's constant use it; in \lean{erdos-770} the fifteen are inherited from \lean{lls-grh}.}\label{tab:compute}
\end{table}

\section{Shared analytic formulas}\label{sec:toolkit}

Most of the analysis rests on a small toolkit, built first for Erd\H{o}s 985 \cite{erdos985}. Put $\xi(s)=s(s-1)\pi^{-s/2}\Gamma(s/2)\zeta(s)$, and let $\rho$ run over its zeros with multiplicity $m_\rho$. For a primitive non-principal character $\chi$ modulo $N$, with parity $a\in\{0,1\}$, put $\Lam(s,\chi)=\pi^{-(s+a)/2}\Gamma(\frac{s+a}2)L(s,\chi)$.

\subsection*{Hadamard's product} Both functions are entire of order one, so for suitable $b$
\begin{equation}\label{eq:hadamard}
\frac{f'}{f}(z)=b+\sum_\rho m_\rho\Bigl(\frac1{z-\rho}+\frac1\rho\Bigr).
\end{equation}
Under RH, combining \eqref{eq:hadamard} at $s$ and at $1-\bar s$ gives an identity that we use repeatedly:
\begin{equation}\label{eq:poisson}
\Re\frac{\xi'}{\xi}(s)=\sum_\rho m_\rho\,\frac{\Re s-\frac12}{|s-\rho|^2}.
\end{equation}
Its right side has a sign on each side of the critical line. At $s=1$ it also yields
\begin{equation}\label{eq:zerosum}
\sum_\rho\frac{m_\rho}{|\rho|^2}=2+\gamma-\log4\pi<\frac1{20}.
\end{equation}
The analogue for $\chi$ is $\sum_\rho m_\rho/|\rho(\rho+1)|\le\frac23(\log N+\frac53)$.

\subsection*{Weighted explicit formulas} Smooth weights make the sums over zeros converge absolutely, so no truncated Perron formula is needed. For primitive $\chi\ne\chi_0$, with the Ces\`aro weight and the kernel $x^s/(s(s+1))$,
\begin{equation}\label{eq:cesaro}
\sum_{n<x}\Lam(n)\chi(n)\Bigl(1-\frac nx\Bigr)=-\sum_\rho m_\rho\frac{x^\rho}{\rho(\rho+1)}+O_\chi(1+\log x).
\end{equation}
With the weight $\log(x/n)$ and the kernel $x^s/s^2$,
\begin{equation}\label{eq:logweight}
\sum_{n\le x}\Lam(n)\chi(n)\log\frac xn=-\sum_\rho m_\rho\frac{x^\rho-1}{\rho^2}-\frac{\Lam'}{\Lam}(0,\chi)\log x+\widetilde\Gamma_\chi(x).
\end{equation}
Under GRH, \eqref{eq:cesaro} gives the estimates of Bach type that McGown, Trevi\~no and Trudgian state as their Lemma~1 without proof. For an integer $x\ge1$, a prime $p$ and $\chi\ne\chi_0$ modulo $p$,
\begin{equation}\label{eq:bach}
\Bigl|\sum_{n<x}\Lam(n)\Bigl(1-\frac nx\Bigr)-\frac x2\Bigr|\le\frac{\sqrt x}{20}+3,\qquad
\Bigl|\sum_{n<x}\Lam(n)\chi(n)\Bigl(1-\frac nx\Bigr)\Bigr|\le c(p,x)\sqrt x\log p,
\end{equation}
where
\[
c(p,x)=\frac23\Bigl(1+\frac2{\sqrt x}+\frac3{x^{3/2}}\Bigr)\Bigl(1+\frac5{3\log p}\Bigr)+\frac{\log x+2}{\sqrt x\log p}.
\]
The first bound in integrated form is $|\Psi_1(t)-t^2/2|\le t^{3/2}/20+4t$, with $\Psi_1(t)=\int_0^t\psi$. Full proofs of \eqref{eq:cesaro} and \eqref{eq:bach} are in \cite{erdos985}.

\section{Named conjectures}\label{sec:named}

\begin{theorem}[Erd\H{o}s 985, \repo{erdos-985}]\label{thm:985}
Every odd prime $p$ has a prime primitive root $q<p$.
\end{theorem}
Let $\omega=\omega(p-1)$. If no prime $q<p$ is a primitive root, expanding the product
\[
\prod_{\ell\mid p-1}\Bigl(\bigl(1-\tfrac1\ell\bigr)-\tfrac1\ell\sum_{0<i<\ell}\eta(n)^{i(p-1)/\ell}\Bigr),\qquad \ord\eta=p-1,
\]
which vanishes at every $n$ that is not a primitive root, gives
\begin{equation}\label{eq:indicator}
S_0\le(2^{\omega}-1)\max_{\chi\ne\chi_0}|S_\chi|,
\end{equation}
where $S_\chi$ is the sum in \eqref{eq:bach} at $x=p$. By \eqref{eq:bach} this fails for $p\ge10^8$. The smaller primes are checked.

\begin{theorem}[Lindel\"of, \repo{lindelof-rh}]
For every $\eps>0$, $\zeta(\frac12+it)=O(t^\eps)$, and $L(\frac12+it,\chi)=O_{\chi,\eps}(t^\eps)$ for every Dirichlet character $\chi$.
\end{theorem}
By \eqref{eq:poisson}, $|\xi(\sigma+it)|$ increases in $\sigma\ge\frac12$, so for $\delta>0$
\[
|\zeta(\tfrac12+it)|\ll t^{\delta/2}\,|\zeta(\tfrac12+\delta+it)|.
\]
Borel and Carath\'eodory's inequality followed by Hadamard's three lines theorem bounds $\log|\zeta(\frac12+\delta+it)|$ by $C(\log t)^{\alpha}$ with $\alpha<1$. This monotonicity replaces the convexity principle and Stirling's formula in Littlewood's proof.

\begin{theorem}[Density hypothesis, \repo{density-rh}]
For every $\eps>0$ there is $C$ such that, for $\frac12\le\sigma\le1$ and $T\ge2$, the number $N(\sigma,T)$ of zeros with $\Re\rho\ge\sigma$ and $0<\Im\rho\le T$ satisfies $N(\sigma,T)\le C\,T^{2(1-\sigma)+\eps}$.
\end{theorem}
Under RH the count is zero for $\sigma>\frac12$, so only $N(\frac12,T)\ll T\log T$ is needed. The repository proves that bound with no hypothesis, from Jensen's inequality.

\begin{theorem}[Grosswald, \repo{grosswald-grh}]
For every prime $p>409$ the least primitive root satisfies $g(p)<\sqrt p-2$.
\end{theorem}
The argument of Theorem~\ref{thm:985} is run at $x=\lfloor\sqrt p\rfloor-2$, with a sieve. If $S$ is a set of primes dividing $p-1$ and $\delta=1-\sum_{\ell\in S}1/\ell>0$, a counterexample satisfies
\[
S_0\le\Bigl(2^{\omega-|S|}\Bigl(2+\frac{|S|-1}{\delta}\Bigr)-1\Bigr)c(p,x)\sqrt x\log p .
\]
This settles $p\ge10^{47}$ with $S=\emptyset$ (where the factor is $2^\omega-1$). A tree search over the prime factors of $p-1$ brings the threshold down to $10^{13}$, and every prime from $409$ to there is checked directly.

\begin{theorem}[Idoneal numbers, \repo{weinberger-grh}]
A positive integer is not of the form $ab+bc+ca$ with $0<a<b<c$ if and only if it is one of Euler's 65 numbers, the largest being $1848$.
\end{theorem}
Since $n=ab+bc+ca$ exactly when $n+a^2=(a+b)(a+c)$, an idoneal $n$ forces $-n$ to be a non-residue modulo every odd prime $p\nmid n$ with $p^2<n$. Let $\psi$ be the primitive character attached to $(-4n/\cdot)$ and $x=\lfloor\sqrt n\rfloor$. Then
\[
\sum_{m<x}\Lam(m)\Bigl(1-\frac mx\Bigr)\bigl(1+\Re\psi(m)\bigr)\le2\,\#P\,\log x,
\]
where $P$ is the set of primes $p<x$ with $p^2<x$ or $p\mid2n$. GRH makes the left side close to $x/2$, a contradiction for $n\ge10^8$. The route avoids genus theory and the class number formula.

\begin{theorem}[ERH for abelian fields, \repo{abelian-erh}]
If $K/\Q$ is a finite abelian extension, every zero of $\zeta_K$ with $0<\Re s<1$ has $\Re s=\frac12$.
\end{theorem}
By the Kronecker-Weber theorem, $K\subseteq\Q(\zeta_N)$ for some $N$. Let $X(K)$ be the characters modulo $N$ trivial on the subgroup fixing $K$. Comparing logarithms gives
\begin{equation}\label{eq:dedekind}
\zeta_K(s)=E(s)\prod_{\chi\in X(K)}L(s,\chi),
\end{equation}
with $E$ holomorphic and zero-free on $\Re s>0$; ramified primes go into $E$ and need no separate treatment. Kronecker-Weber and the continuation of $\zeta_K$ come from \cite{rhrepo}.

\section{Statements equivalent to RH}\label{sec:equiv}

\begin{theorem}[Robin, \repo{robin-rh}]
$\sigma(n)<e^\gamma n\log\log n$ for every $n>5040$, and $\sigma(n)\le H_n+e^{H_n}\log H_n$ for every $n\ge1$.
\end{theorem}
Both ranges start from the envelope that defines colossally abundant numbers: for every $\eps>0$,
\begin{equation}\label{eq:robin}
\log\frac{\sigma(n)}n\le\eps\log n+\sum_p\max_{k\ge0}\Bigl(\log\frac{\sigma(p^k)}{p^k}-k\eps\log p\Bigr).
\end{equation}
For $\log n\ge e^{30}$ one takes $\eps=\log(1+1/x)/\log x$ with $x=\log n$ and uses $\Psi_1$. Below that the right side of \eqref{eq:robin} is a line in $x$ and $\gamma+\log\log x$ is concave, so checking the two ends of an interval settles every $n$ in between.

\begin{theorem}[Nicolas, \repo{nicolas-rh}]
With $N_k=p_1\cdots p_k$, $N_k/\varphi(N_k)>e^\gamma\log\log N_k$ for every $k\ge1$.
\end{theorem}
Let $S(x)=-\sum_{p\le x}\log(1-1/p)$, $w(t)=(1+\log t)/(t^2\log^2t)$ and $T(x)=\sum_{p>x}(-\log(1-1/p)-1/p)$. Then
\[
S(x)-\gamma-\log\log\theta(x)\ \ge\ \int_x^\infty(t-\theta(t))\,w(t)\,dt-T(x).
\]
The gain is $\psi-\theta\ge\psi(\sqrt t)$, which outweighs the zero sum for $x\ge10^5$.

\begin{theorem}[Li, \repo{li-criterion-rh}]
$\lambda_n=\sum_\rho\bigl(1-(1-1/\rho)^n\bigr)>0$ for every $n\ge1$.
\end{theorem}
Here zeros are paired by conjugation and, as in the \emph{formal-conjectures} statement, counted once each. On the critical line $w=1-1/\rho$ has $|w|=1$, so
\begin{equation}\label{eq:li}
2\Re(1-w^n)=|1-w^n|^2\ge0,\qquad |1-w^n|\le\frac n{|\rho|}.
\end{equation}
The sum converges by \eqref{eq:zerosum}. It is positive because $\zeta$ has infinitely many zeros, which the repository proves without Stirling's formula.

\begin{theorem}[Speiser, \repo{speiser-rh}]
$\zeta'(s)\ne0$ for $0<\Re s<\frac12$.
\end{theorem}
If $\zeta'(s)=0$ in the strip, then
\begin{equation}\label{eq:speiser}
\frac{\xi'}\xi(s)=\frac1s+\frac1{s-1}-\frac12\log\pi+\frac12\,\frac{\Gamma'}\Gamma\Bigl(\frac s2\Bigr).
\end{equation}
For $|\Im s|\ge7$ the right side has positive real part while the left is at most $0$ by \eqref{eq:poisson}. For $|\Im s|<\frac52$ the two sides have incompatible sizes. The band between is covered by discs on which an Euler-Maclaurin expansion keeps $\zeta'$ away from zero.

\begin{theorem}[Mertens bound, \repo{mertens-rh}]
For every $\eps>0$, $M(x)=\sum_{n\le x}\mu(n)=O(x^{1/2+\eps})$. More precisely, for all large $x$,
\[
|M(x)|\le\sqrt x\,\exp\Bigl(\frac{25\log x}{\log\log x}\Bigr).
\]
\end{theorem}
The second bound is Titchmarsh's \cite{titchmarsh27} with an explicit constant. With $A(N)=\sum_{n<N}\mu(n)(N-n)$ one has $M(N)=A(N+1)-A(N)$ and
\[
A(N)=\frac1{2\pi i}\int_{(\sigma_0)}\frac{N^{s+1}}{s(s+1)\,\zeta(s)}\,ds,\qquad\sigma_0>\tfrac12,
\]
so a bound for $1/\zeta$ to the right of the critical line is enough. The repository also derives the bounds of Riesz, Hardy and Littlewood, Franel and Landau, the Liouville function, and $\det R_n=M(n)$ for Redheffer's matrix.

\section{Further consequences}\label{sec:further}

\begin{theorem}[Lamzouri, Li and Soundararajan, \repo{lls-grh}]
\begin{enumerate}
\item For every prime $q\ge5$ there is a quadratic non-residue $n\le(\log q)^2$.
\item For primitive $\chi$ modulo $q\ge10^{10}$,
\[
|L(1,\chi)|\le2e^\gamma\Bigl(\log\log q-\log2+\frac12+\frac1{\log\log q}\Bigr).
\]
\item Every proper subgroup of $(\Z/q\Z)^\times$, $q\ge3$, omits a unit $n\le2(\log q)^2$.
\end{enumerate}
\end{theorem}
Parts (1) and (2) are from \cite{lls}; part (3) is Bach's theorem \cite{bach90} with his constant. The repository covers every theorem of \cite{lls}, including least primes in cosets and class number bounds. The main tool is \eqref{eq:logweight}, with the zero sum evaluated through
\begin{equation}\label{eq:z2}
\sum_\rho\frac{m_\rho}{|\rho|^2}=\log N+\Re\frac{\Lam'}\Lam(1,\chi)+\Re\frac{\Lam'}\Lam(1,\bar\chi).
\end{equation}

\begin{theorem}[Bach and Sorenson, \repo{bach-sorenson-grh}]
If $q\ge2$ and $(a,q)=1$, some prime $p\equiv a\pmod q$ satisfies $p<2(q\log q)^2$.
\end{theorem}
Bounding each character separately stops at the constant $\frac72$. Summing over all characters first, with
\[
\sum_{\chi\bmod q}\Re\frac{L'}L(\sigma,\chi)\le0\qquad(\sigma>1),
\]
gives the constant $2$ for $q\ge2^{15}$.

\begin{theorem}[Schoenfeld, \repo{schoenfeld-rh}]
\begin{align*}
|\psi(x)-x|&<\frac{\sqrt x\log^2x}{8\pi}\quad(x\ge73.2),&
|\theta(x)-x|&<\frac{\sqrt x\log^2x}{8\pi}\quad(x\ge599),\\
|\pi(x)-\li(x)|&<\frac{\sqrt x\log x}{8\pi}\quad(x\ge2657).
\end{align*}
\end{theorem}
Differencing $\Psi_1$ over a step $h\asymp\sqrt x$ gives
\[
|\psi(x)-x|\le\frac h2+O(1)+\sqrt{x+h}\;F\Bigl(\frac{2(x+h)}h\Bigr),\qquad F(T)=\sum_\rho m_\rho\min\Bigl(\frac1{|\rho|},\frac T{|\rho|^2}\Bigr).
\]
Rosser's explicit zero count then bounds $F(T)$ by $\log^2T/(2\pi)$ up to explicit lower-order terms. The formal proof of Rosser's count is vendored from public Lean projects credited in the repository.

\begin{theorem}[Cram\'er; Carneiro, Milinovich and Soundararajan, \repo{cramer-rh}]
For $x\ge4$ the interval $[x,\,x+\frac{22}{25}\sqrt x\log x]$ contains a prime. Moreover
\[
\int_1^X\Bigl(\frac{\psi(x)-x}x\Bigr)^2dx\sim c_1\log X,\qquad c_1=\sum_{\rho}\frac{m_\rho^2}{|\rho|^2}\quad(\rho\text{ distinct}).
\]
\end{theorem}
The gap bound uses a piecewise linear weight $w(n)=\sum_ja_j\max(n_j-n,0)$, positive on the gap and negative on either side. The explicit formula reads
\[
\sum_n\Lam(n)w(n)=\sum_ja_j\Psi_1(n_j),
\]
and the primes outside the gap are bounded by the large sieve form of Brun and Titchmarsh's inequality.

\begin{theorem}[Erd\H{o}s 770(iii), \repo{erdos-770}]\label{thm:770}
Let $h(n)$ be the least $M>2$ such that $2^n-1,\dots,M^n-1$ have greatest common divisor $1$, and let $P(n)$ be the largest prime $p$ with $p-1\mid n$. For every $n\ge2$,
\[
P(n)\le h(n)\le\max\bigl(P(n),\lceil200\log^2n\rceil\bigr).
\]
In particular, for every $\eps>0$ and all large $n$, $P(n)>n^\eps$ implies $h(n)=P(n)$.
\end{theorem}
If a prime $r$ divides $i^n-1$ for all $2\le i\le M$, the kernel of $u\mapsto u^n$ on $(\Z/r\Z)^\times$ contains $1,\dots,M$. When $r\le n^{10}$ and the kernel is proper, part (3) of the theorem above gives $M<2\log^2r\le200\log^2n$. When $r>n^{10}$, products of small primes give more than $n$ elements of a group of order at most $n$. That GRH implies a bound of this shape is remarked without proof in notes of Korsky, which we know through secondary sources.

\begin{theorem}[Selberg, \repo{erdos-233}]
Let $d_n=p_{n+1}-p_n$. Then $\sum_{n<N}d_n^2=O(N\log^4N)$. If $h(x)/\log^2x\to\infty$, almost all intervals $(x,x+h(x)]$ contain a prime; if also $h$ is increasing and $h(x)=o(x)$, then $\pi(x+h(x))-\pi(x)\sim h(x)/\log x$ for almost all $x$.
\end{theorem}
Put $\Delta_\delta(y)=\Psi_1(y(1+2\delta))-2\Psi_1(y(1+\delta))+\Psi_1(y)$. Its mean square satisfies
\[
\int_x^{4x}\bigl(\Delta_\delta(y)-\delta^2y^2\bigr)^2\,dy\ll x^4\delta^3\log^2\frac2\delta+x,
\]
and a long gap forces $\Delta_\delta(y)$ to be far from $\delta^2y^2$ on a set of comparable length. The conjectured bound $O(N\log^2N)$ of Problem 233 is open even under RH.

\section{Method}\label{sec:method}

\subsection*{Moves taken from mathematicians} We began by reading what eleven mathematicians said about how they work (Tao, Grothendieck, Erd\H{o}s, Gowers, Wiles, Thurston, Zeilberger, Ramanujan, Scholze, Arnold and Perelman). Each account was reduced to moves an agent can carry out: isolate one difficulty at a time, try to prove and to disprove at once, repair a false guess, look for the one missing estimate. Where the record was thin, as for Ramanujan and Perelman, we did not invent a process. The moves became roles. One agent reduces a goal to named lemmas, others search for counterexamples, try weaker statements or read the literature, and one writes Lean. A dependency graph records the open lemmas, and when one fails the results above it are reopened. The first version is \repo{math-superagent}.

\subsection*{An atlas of problems} Agents asked to solve an open problem tend to choose a famous one. We therefore compiled 2{,}196 problems, 1{,}409 of them unsolved, from \emph{Erd\H{o}s Problems} \cite{erdosproblems}, \emph{formal-conjectures}, Wikipedia and about a hundred entries of our own; 808 of the unsolved ones have a Lean statement. Each problem has a difficulty score for mathematicians and a score for how well it suits agents, and problems are ranked by
\[
\text{opportunity}=\text{agent fit}-\text{difficulty}.
\]
As a retrospective test we scored solved problems on what was known before their proofs. Those settled between 2019 and 2024, such as Ringel's conjecture and Kaplansky's unit conjecture, came out positive; Fermat's Last Theorem and the Poincar\'e conjecture came out negative.

\subsection*{The loop} RSI Ralph repeats a prompt and carries forward only the files an agent leaves. It runs in epochs, tests each proposed change against the current best on the same inputs, and may edit the team's own prompts, roles and tools. The task, the team and the fixed rules are three JSON files, so a change is a small diff that can be scored and rolled back. Agents cannot edit the scorer. For a formal proof the score is a ladder of Lean statements, and a change is kept only if the project builds and \lean{\#print axioms} shows nothing outside the allowed list. A lead session rebuilt each branch itself before merging it. The first runs on Erd\H{o}s problems used \href{https://github.com/tinyhumansai/supermath}{\texttt{supermath}}.

\subsection*{Memory} Papers, proof blueprints, harness logs and the atlas were stored in Living Brain, which compiles captured material into linked, versioned pages that keep their sources. When the proof of RH appeared, the question was which problems had been waiting for exactly \eqref{eq:rh} and \eqref{eq:grh}. Recall over this material surfaced the candidates and recorded why others failed. Artin's primitive root conjecture needs GRH for Kummer fields, which \eqref{eq:grh} does not include; binary Goldbach and the twin prime conjecture are not known to follow from RH.

Living Brain makes four decisions about each item: whether it deserves upkeep, where it belongs, what it should be linked to, and when it should return. The last of these matters most in this work, since a proof is resumed many times. Separate health measures track whether the store is fresh, ready and under strain. They rest on a decay for an event of age $a$ with half-life $h$ and a saturating scale for an aggregate $x$,
\[
D(a,h)=2^{-a/h},\qquad S(x,k)=100\bigl(1-e^{-\max(0,x)/k}\bigr).
\]
None of these numbers measures whether a proof is correct.

\subsection*{The recipe for a target} Each project followed the same steps.
\begin{enumerate}
\item Write the target statement and fix it.
\item Prove $\mathrm{RH}\Rightarrow X$ or $\mathrm{GRH}\Rightarrow X$, with the hypothesis as an argument, by splitting the proof into lemmas for separate agents.
\item Sharpen the constants until the threshold is within reach of a computation, then run a checker that has been proved sound.
\item Supply the hypothesis from \cite{rhrepo}.
\end{enumerate}
Erd\H{o}s 985 came first, on 22 November, and its toolkit (Section~\ref{sec:toolkit}) was reused everywhere else. Classical theorems missing from Mathlib, such as Kronecker-Weber and the class number formula, entered as named hypotheses and were discharged from \cite{rhrepo} in the last step.

\section{Remarks}

\subsection*{Constants} A first proof often landed above the published constant. Bach's bound went from $\frac{21}5$ to $\frac{11}4$ to $2$, the bound of Bach and Sorenson from $\frac72$ to $2$, Schoenfeld's from $\frac18$ to $\frac1{8\pi}$, and Cram\'er's gap constant from $2$ to $\frac{22}{25}$. Each time the obstacle was one lossy step, and keeping that step exact removed it.

\subsection*{What did not work} A search of the open Erd\H{o}s problems for further consequences of RH by a new argument found nothing we kept. Problem 233 asks for more than RH is known to give.

\subsection*{Limits of the evidence} The proofs show that the statements hold. They do not show how much the memory contributed. Recall was used to choose targets, but we ran no comparison without it. A fair test would give the same model, roles and budget to runs with and without memory, and compare how well each ranks an expert list of consequences, the time to a first accepted proof, repeated dead ends, and cost. Our own checks link against the community build of \cite{rhrepo} at commit \lean{507a75d91d} and do not repeat the kernel checks of \cite{sutherland}.

\subsection*{Acknowledgments and disclosure} We thank Levent Alp\"oge, the Lean and Mathlib communities, and the authors of the vendored Lean projects named in each repository. The formal proofs were written by AI agents, using GLM models by Z.ai served on GMI Cloud, with literature search by Exa. Steven Enamakel is the founder of OpenHuman and Christel Buchanan is the founder of Living Brain, the systems described in Section~\ref{sec:method}.

\begin{thebibliography}{99}
\bibitem{alpoge} L. Alp\"oge, Lean proof of the Riemann hypothesis, announced 20 November 2026, \url{https://alpo.ge/remark.pdf}.
\bibitem{bach90} E. Bach, \emph{Explicit bounds for primality testing and related problems}, Math. Comp. \textbf{55} (1990), 355-380.
\bibitem{bachsorenson} E. Bach and J. Sorenson, \emph{Explicit bounds for primes in residue classes}, Math. Comp. \textbf{65} (1996), 1717-1735.
\bibitem{cms} E. Carneiro, M. Milinovich and K. Soundararajan, \emph{Fourier optimization and prime gaps}, Comment. Math. Helv. \textbf{94} (2019), 533-568; arXiv:1708.04122.
\bibitem{erdos985} S. Enamakel, \emph{Every odd prime has a smaller prime primitive root under GRH: a formalized proof of Erd\H{o}s Problem 985}, 2026, \url{https://github.com/senamakel/erdos-985/tree/main/paper}.
\bibitem{erdosproblems} T. Bloom, \emph{Erd\H{o}s Problems}, \url{https://www.erdosproblems.com}.
\bibitem{formalconj} Google DeepMind, \emph{formal-conjectures}, \url{https://github.com/google-deepmind/formal-conjectures}.
\bibitem{lls} Y. Lamzouri, X. Li and K. Soundararajan, \emph{Conditional bounds for the least quadratic non-residue and related problems}, Math. Comp. \textbf{84} (2015), 2391-2412; arXiv:1309.3595.
\bibitem{leanfro} Lean FRO and Mathlib Initiative, \emph{Lean verification of the Riemann Hypothesis}, 21 November 2026, \url{https://leanprover-community.github.io/blog/posts/riemann-hypothesis/}.
\bibitem{li} X.-J. Li, \emph{The positivity of a sequence of numbers and the Riemann hypothesis}, J. Number Theory \textbf{65} (1997), 325-333.
\bibitem{littlewood} J. E. Littlewood, \emph{Quelques cons\'equences de l'hypoth\`ese que la fonction $\zeta(s)$ de Riemann n'a pas de z\'eros dans le demi-plan $\Re(s)>\frac12$}, C. R. Acad. Sci. Paris \textbf{154} (1912), 263-266.
\bibitem{livingbrain} \emph{Living Brain}, \url{https://livingbrain.com}.
\bibitem{mtt} K. McGown, E. Trevi\~no and T. Trudgian, \emph{Resolving Grosswald's conjecture on GRH}, Funct. Approx. Comment. Math. \textbf{55} (2016), 215-225; arXiv:1508.05182.
\bibitem{nicolas} J.-L. Nicolas, \emph{Petites valeurs de la fonction d'Euler}, J. Number Theory \textbf{17} (1983), 375-388.
\bibitem{openhuman} \emph{OpenHuman}, \url{https://github.com/tinyhumansai/openhuman}.
\bibitem{rhrepo} \emph{leanprover-community/riemann-hypothesis}, \url{https://github.com/leanprover-community/riemann-hypothesis}.
\bibitem{robin} G. Robin, \emph{Grandes valeurs de la fonction somme des diviseurs et hypoth\`ese de Riemann}, J. Math. Pures Appl. \textbf{63} (1984), 187-213.
\bibitem{rsiralph} \emph{RSI Ralph}, \url{https://github.com/tinyhumansai/rsi-ralph}.
\bibitem{schoenfeld} L. Schoenfeld, \emph{Sharper bounds for the Chebyshev functions $\theta(x)$ and $\psi(x)$. II}, Math. Comp. \textbf{30} (1976), 337-360.
\bibitem{selberg} A. Selberg, \emph{On the normal density of primes in small intervals, and the difference between consecutive primes}, Arch. Math. Naturvid. \textbf{47} (1943), no. 6, 87-105.
\bibitem{speiser} A. Speiser, \emph{Geometrisches zur Riemannschen Zetafunktion}, Math. Ann. \textbf{110} (1934), 514-521.
\bibitem{sutherland} A. Sutherland, \emph{Verification files for the Lean proof of the Riemann Hypothesis}, \url{https://math.mit.edu/~drew/grh/}.
\bibitem{titchmarsh27} E. C. Titchmarsh, \emph{A consequence of the Riemann hypothesis}, J. London Math. Soc. \textbf{2} (1927), 247-254.
\bibitem{washington} L. C. Washington, \emph{Introduction to Cyclotomic Fields}, 2nd ed., Springer, 1997.
\bibitem{weinberger} P. J. Weinberger, \emph{Exponents of the class groups of complex quadratic fields}, Acta Arith. \textbf{22} (1973), 117-124.
\end{thebibliography}
\end{document}
